\section{开源、未来方向与课堂问答}

课堂最后没有把 NLLB 描述成终点，而是提出三条未来方向，并在 Q\&A 中补充社区参与、数据工作量、Common Crawl、zero-shot、speech translation 与 foundation model 的边界。这部分口头内容解释了 2022 年专门化 NLLB 系统如何连接后来的通用多模态模型。

\subsection{Explicit multilinguality、support 与易用性}

未来方向 slide 强调 explicit multilinguality、support for everyone，以及更容易使用和训练。第一点要求模型显式知道语言/脚本与方向；第二点把尾部语言纳入目标；第三点则承认，只有大型机构能运行的数据管线仍不足以形成普惠基础设施。

\lecturefigure{slide-39-future-directions.jpg}{NLLB 的未来方向：显式多语、更广支持与更低使用门槛}{Stanford 课堂视频 00:41:15}

\begin{knowledgebox}{读图：未来方向对应表示、覆盖与工程三个层面}
Explicit multilinguality 作用于 language tags、routing 与 evaluation；support for everyone 作用于尾部语言、方言与社区需求；ease of use/training 作用于开放工具、算力门槛和本地适配。三者缺一，研究结果都难转成持续可用的公共能力。
\end{knowledgebox}

\teachervoice{Fan 特别指出，一些语言主要是口语而非书面语。文本 NLLB 无法凭空覆盖这些语言；未来必须连接 speech、transcription 与 text translation，而不是把“无网页数据”当成语言不存在。}

\subsection{数据工作与模型工作大致同等重要}

Q\&A 中有人询问项目资源分配，Fan 用近似口径说，大约一半是 data 或 data-adjacent work，另一半是 modeling，但边界很模糊。LID、encoder、filtering 和 evaluation 同时属于数据与模型；这个回答的重点不是精确项目管理比例，而是反驳“主要工作是训练大模型”。

\begin{importantbox}{NLLB 的工程资产层级}
\begin{tabularx}{\textwidth}{@{}L{0.2\textwidth}L{0.29\textwidth}X@{}}
\toprule
层级 & 公开资产 & 可复现的问题 \\
\midrule
Need / standards & 访谈结论与语言规范过程 & 目标与标签为何这样定义 \\
Evaluation & FLORES-200、human protocol、Toxicity-200 & 模型是否进步、哪里不安全 \\
Data tools & LID、LASER3、Stopes、重建脚本 & 训练数据如何产生与过滤 \\
Models & NLLB checkpoints 与代码 & 最终推理行为与适配 \\
\bottomrule
\end{tabularx}
完整开放意味着尽量覆盖从问题到模型的链条，而不是只有 checkpoint。
\end{importantbox}

\teachervoice{老师说“50-50”时立即补充 data 与 modeling 很难严格分开。讲义把它作为数量级直觉：大型多语项目的主要创新有相当部分发生在数据、评测和工具链。}

\subsection{Common Crawl 的开放不等于均匀代表}

课堂说明 Common Crawl 是可下载的开放网页快照，并对更新频率保留了不确定措辞。开放访问解决了数据获取许可的一部分问题，却不解决语言网页占比、重复、模板、垃圾内容、地域偏差与时间漂移。

\begin{warningbox}{Web-scale 与 language coverage 是两个不同变量}
一个 corpus 可以有数十亿网页，同时对某些语言几乎没有有效文本。评估 mining pipeline 时应报告每语言 raw pages、LID pass rate、去重后 token、候选 bitext、人工 precision 与领域分布，而不是只报总 crawl 大小。
\end{warningbox}

\subsection{Zero-shot 的边界：未见方向，不等于未见语言}

NLLB 没有为每个 $A\rightarrow B$ 方向挖掘训练数据；组合爆炸使这既昂贵也不一定符合需求。某个方向可以是 zero-shot，但模型在其他方向见过语言 $A$ 和 $B$，因此仍可通过共享 encoder/decoder 表示迁移。

\begin{importantbox}{两种“未见”必须分开}
\begin{itemize}
  \item \textbf{Unseen direction}：训练没有直接出现 $A\rightarrow B$，但出现过 $A\leftrightarrow E$、$B\leftrightarrow E$ 等任务。
  \item \textbf{Unseen language}：语言 $A$ 的文本、token 或监督几乎从未出现。
\end{itemize}
课堂 zero-shot 主要指前者。把它写成对完全新语言的翻译能力，会夸大泛化范围。
\end{importantbox}

\teachervoice{Fan 解释，某些 zero-shot pair 表现很好，是因为模型见过输入语言和输出语言，只是没有见过该有向组合；总体上这类方向仍更弱，而且不是对完全 unseen language 的零样本。}

\subsection{从文本 NLLB 到 SeamlessM4T}

针对语音问题，Fan 明确说 NLLB 本身主要是文本翻译；后续 SeamlessM4T 才把 speech、transcription 与 text translation 联合起来。这个回答为项目边界提供了清晰分层：文本 pipeline 的数据假设不能直接外推到主要依赖口语的语言。

\begin{knowledgebox}{文本与语音扩展的证据链}
Speech translation 需要音频、转写、说话人/口音覆盖、语音分段与声学表示；文本 NLLB 主要处理书面语句对。SeamlessM4T 通过共享多模态表示连接两者，但仍需新的语音评测与安全分析，不能只复用 FLORES 文本分数。
\end{knowledgebox}

\subsection{Foundation model 时代的再解释}

课堂最后讨论，如果在更强基础模型时代重启项目，可能使用大 foundation model 加 supervised fine-tuning，减少为每个任务设计专门系统的需要。Fan 回顾早期 NLP 中 translation、summarization、QA 研究彼此分离，而现在越来越多任务由共享底座驱动。

\teachervoice{老师的判断是方向性而非定理：基础模型越强，领域 fine-tuning 可能越少，但低资源语言是否在预训练中得到足够表示仍是前提。通用底座不能自动修复训练语料中的语言不平衡。}

\begin{warningbox}{不要用后来的 LLM 叙事抹掉 NLLB 的贡献}
更强 foundation model 仍需要多语数据、可靠 benchmark、language identification、人工评测和 safety probes。NLLB 的长期价值恰恰是把这些基础设施与模型联合公开，而不是只给出一种 2022 年架构。
\end{warningbox}

\subsection{本章小结}

未来多语系统要覆盖显式语言表示、尾部社区和更低工程门槛。Q\&A 进一步限定：数据与模型同等关键，开放网页并不均衡，zero-shot 多指未见方向，speech 需要新模态证据，foundation model 也不能替代多语基础设施。

